
\documentclass[12pt]{article}

\usepackage{geometry}
\usepackage{titling}
\usepackage{amsmath,amssymb}
\usepackage{fontspec}
\usepackage{unicode-math}
\usepackage{array}
\usepackage{colortbl}
\usepackage[dvipsnames]{xcolor}
\usepackage[most]{tcolorbox}

\geometry{a4paper, top=2.5cm, bottom=2.5cm, left=1.5cm, right=1.5cm}
\setlength{\droptitle}{-3cm}

\setmonofont{DejaVu Sans Mono}
\setmainfont{Libertinus Serif}
\setsansfont{Libertinus Sans}
\setmathfont{STIXTwoMath-Regular.otf}

\definecolor{question}{rgb}{0,0.1,1.0}

\newcommand{\vsp}{\vspace*{1em}}
\newcommand{\QuizQuestion}[1]{{\texts\textbf{\color{question}{#1}}}}

\title{Linear Functions - Matrix-Vector Multiplication}
\author{Devi Prasad M}
\date{August 2026}

\begin{document}
\pagecolor{yellow!10}

\maketitle

\section{Matrix-Vector Multiplication}


Let $\mathbf{A} \in \mathbb{R}^{m \times n}$ be a matrix,
$\mathbf{u}, \mathbf{v} \in \mathbb{R}^n$ be vectors,
and $c \in \mathbb{R}$, a scalar.

\vsp

\noindent Matrix-vector multiplication is linear because which satisfies
\begin{enumerate}
    \item \textbf{Additivity:}
    \[
        \mathbf{A}(\mathbf{u} + \mathbf{v}) = \mathbf{A} \mathbf{u} + \mathbf{A} \mathbf{v}.
    \]

    \item \textbf{Homogeneity:}
    \[
        \mathbf{A}(c\mathbf{u}) = c(\mathbf{A} \mathbf{u}).
    \]
\end{enumerate}

\noindent Both properties follow directly from the definition of matrix multiplication and the distributive and associative properties of real numbers.
Hence, the function $\mathbf{x} \mapsto \mathbf{A}\mathbf{x}$ is linear.

\vsp

\noindent Let us assume representative values for $\mathbf{A}$, $\mathbf{u}$, and $\mathbf{v}$:

\[
    \mathbf{A} = \begin{bmatrix}
        w_{11} & w_{12} & w_{13} \\
        w_{21} & w_{22} & w_{23} \\
        w_{31} & w_{32} & w_{33} \\
        w_{41} & w_{42} & w_{43}
    \end{bmatrix}
    ,\ \ \
    \mathbf{u} = \begin{bmatrix}
        u_1 \\
        u_2 \\
        u_3 \\
    \end{bmatrix}
    ,\ \, \mathrm{and} \ \ \
    \mathbf{v} = \begin{bmatrix}
        v_1 \\
        v_2 \\
        v_3 \\
    \end{bmatrix}
\]

\subsection{Additivity}

\subsubsection*{\noindent Compute $\mathbf{A} (\mathbf{u} + \mathbf{v})$:}

\[
    \begin{bmatrix}
        w_{11} & w_{12} & w_{13} \\
        w_{21} & w_{22} & w_{23} \\
        w_{31} & w_{32} & w_{33} \\
        w_{41} & w_{42} & w_{43}
    \end{bmatrix}
    \,
    \begin{bmatrix}
        (u_1 + v_1) \\
        (u_2 + v_2) \\
        (u_3 + v_3) \\
    \end{bmatrix}
    =
    \begin{bmatrix}
        w_{11} (u_1 + v_1)  + w_{12} (u_2 + v_2) + w_{13} (u_3 + v_3) \\
        w_{21} (u_1 + v_1)  + w_{22} (u_2 + v_2) + w_{23} (u_3 + v_3) \\
        w_{31} (u_1 + v_1)  + w_{32} (u_2 + v_2) + w_{33} (u_3 + v_3) \\
        w_{41} (u_1 + v_1)  + w_{42} (u_2 + v_2) + w_{43} (u_3 + v_3)
    \end{bmatrix}
\]

\vsp

\subsubsection*{Compute $\mathbf{A} \mathbf{u}$:}

\[
    \begin{bmatrix}
        w_{11} & w_{12} & w_{13} \\
        w_{21} & w_{22} & w_{23} \\
        w_{31} & w_{32} & w_{33} \\
        w_{41} & w_{42} & w_{43}
    \end{bmatrix}
    \,
    \begin{bmatrix}
        u_1 \\
        u_2 \\
        u_3
    \end{bmatrix}
    =
    \begin{bmatrix}
        w_{11} u_1  + w_{12} u_2 + w_{13} u_3 \\
        w_{21} u_1  + w_{22} u_2 + w_{23} u_3 \\
        w_{31} u_1  + w_{32} u_2 + w_{33} u_3 \\
        w_{41} u_1  + w_{42} u_2 + w_{43} u_3
    \end{bmatrix}
\]

\vsp

\subsubsection*{\noindent Compute $\mathbf{A} \mathbf{v}$:}

\[
    \begin{bmatrix}
        w_{11} & w_{12} & w_{13} \\
        w_{21} & w_{22} & w_{23} \\
        w_{31} & w_{32} & w_{33} \\
        w_{41} & w_{42} & w_{43}
    \end{bmatrix}
    \,
    \begin{bmatrix}
        v_1 \\
        v_2 \\
        v_3
    \end{bmatrix}
    =
    \begin{bmatrix}
        w_{11} v_1  + w_{12} v_2 + w_{13} v_3 \\
        w_{21} v_1  + w_{22} v_2 + w_{23} v_3 \\
        w_{31} v_1  + w_{32} v_2 + w_{33} v_3 \\
        w_{41} v_1  + w_{42} v_2 + w_{43} v_3
    \end{bmatrix}
\]


\vsp

\noindent Compute $\mathbf{A} \mathbf{u} + \mathbf{A} \mathbf{v}$:
\[
    \begin{bmatrix}
        w_{11} u_1  + w_{12} u_2 + w_{13} u_3 \\
        w_{21} u_1  + w_{22} u_2 + w_{23} u_3 \\
        w_{31} u_1  + w_{32} u_2 + w_{33} u_3 \\
        w_{41} u_1  + w_{42} u_2 + w_{43} u_3
    \end{bmatrix}
    + \,
    \begin{bmatrix}
        w_{11} v_1  + w_{12} v_2 + w_{13} v_3 \\
        w_{21} v_1  + w_{22} v_2 + w_{23} v_3 \\
        w_{31} v_1  + w_{32} v_2 + w_{33} v_3 \\
        w_{41} v_1  + w_{42} v_2 + w_{43} v_3
    \end{bmatrix}
    =
    \begin{bmatrix}
        w_{11} (u_1 + v_1)  + w_{12} (u_2 + v_2) + w_{13} (u_3 + v_3) \\
        w_{21} (u_1 + v_1)  + w_{22} (u_2 + v_2) + w_{23} (u_3 + v_3) \\
        w_{31} (u_1 + v_1)  + w_{32} (u_2 + v_2) + w_{33} (u_3 + v_3) \\
        w_{41} (u_1 + v_1)  + w_{42} (u_2 + v_2) + w_{43} (u_3 + v_3)
    \end{bmatrix}
\]

\noindent Note that $\mathbf{A}(\mathbf{u} + \mathbf{v}) = \mathbf{A} \mathbf{u} + \mathbf{A} \mathbf{v}$.
Thus \textbf{Additivity} holds for general matrix-vector multiplication.

\newpage
\pagecolor{yellow!30}
\subsection{Homogeneity}

\begin{center}
\vfill

{\fontsize{50pt}{100pt}\selectfont {\color{blue}\textbf{Prove Homogeneity!}}}
\end{center}

\vfill

\end{document}
